\section{Discussion \& Future Work}\label{sec:discussion}

\section{Discussion : Chemist agent possible improvement}
Trough the whole filtering and scoring process by the agent, the computed scores are obtained using a peptide sequence code (e.g ACDEF). 
From a certain sequence we can obtain multiple other sequences in the SMILES format (e.g ). Therefore we might be losing possible data along the process.
Our lack of chemistery knowledge doesn't allows us to be sure of how it would change the outcome of our results or even the biologist agent prediction. 



\subsection{Future Work: Peptide-Specific Biological Scoring}\label{subsec:future_biologist}
The main improvement is to replace the current proxy with a model trained (or fine-tuned) specifically on peptides. Ideally, this model should score or classify a peptide conditionally on a target (or target-related context), instead of relying on a general protein-space similarity signal.

The expected benefit is a more accurate and reliable biological signal for ranking candidates in the loop. A peptide-specialized model could also provide richer outputs than a single proxy score, for example clearer estimates of expected activity and additional information about risk or toxicity.

In short, our current biologist agent is a practical baseline for iterative optimization, while a target-conditioned peptide model is the natural next step to improve biological relevance and decision quality.
